% ESLab Term Project Proposal -- compile with XeLaTeX
\documentclass[12pt]{extarticle}
\usepackage[a4paper, margin = 2.5cm]{geometry}%A4 paper size, 2.5cm margin
\usepackage{xeCJK}%For Chinese typesetting
\usepackage{indentfirst}%Indent the first paragraph of each section
\usepackage{fontspec}%For font selection
\setmainfont{Times New Roman}%Set main font to Times New Roman
\setCJKmainfont[AutoFakeBold=4, AutoFakeSlant=0.2]{Kaiti TC}%Set main CJK font to Kaiti TC with fake bold and slant
\setCJKmonofont{Kaiti TC}%Provide a CJK fallback for monospaced contexts
\usepackage{amsmath, amssymb, graphicx, float, caption, hyperref, bookmark, booktabs, fancyhdr}%Various useful packages
\hypersetup{colorlinks=true, allcolors=black, pdfstartview=Fit, breaklinks=true}%Hyperlink settings
\usepackage{titlesec}%For section title formatting
\usepackage{tabularx, array, enumitem, xurl}%Tables, lists, URL line breaks
\usepackage{tikz}%System diagram
\usetikzlibrary{positioning, arrows.meta, fit, backgrounds, calc}
\pagestyle{fancy}
\setlength{\headheight}{15pt}
\rhead{}
\lhead{}
\chead{}
\renewcommand{\headrulewidth}{0.4pt}
\renewcommand{\footrulewidth}{0.4pt}
\lhead{嵌入式系統實驗}
\rhead{期末專題提案}
\cfoot{\thepage}
\captionsetup{font=small, labelformat=empty}
\setlist{itemsep=0.1em, topsep=0.3em}
\renewcommand{\arraystretch}{1.25}
\newcolumntype{Y}{>{\raggedright\arraybackslash}X}

\begin{document}

\begin{center}
{\LARGE\bfseries AirTaiko}\\[0.4em]
{\large 以兩塊 STM32 IoT Node 作為無線鼓棒的太鼓節奏遊戲}\\[0.6em]
{\small 電工實驗（嵌入式系統）期末專題提案}
\end{center}

\section{題目、GitHub Repo 與組員}

兩手各拿一塊 STM32 IoT node 當鼓棒，在空中敲打。板子即時偵測敲擊並判斷是「咚」（鼓面）還是「咔」（鼓邊），
透過 BLE 把打擊事件傳到電腦；電腦上執行我們自製的太鼓節奏遊戲，音符隨音樂捲動，依打擊時間判定「良／可／不可」並計分。

\noindent{\small\begin{tabularx}{\linewidth}{@{}lY@{}}
\toprule
題目 & AirTaiko：以兩塊 STM32 IoT Node 作為無線鼓棒的太鼓節奏遊戲 \\
GitHub repo & \url{https://github.com/leezhi1019/ESLab-FinalProject} \\
組員 & 李健誌、黃琮堡、莊景皓 \\
開發板 & STM32 B-L475E-IOT01A × 2（左右鼓棒，必要） \\
無線通訊 & BLE（兩塊 STM32 皆為 GATT server，電腦為 GATT client 同時連線） \\
遊戲端 & 電腦上以 Python（pygame）自製遊戲畫面 \\
\bottomrule
\end{tabularx}}

\section{動機}

本專題延伸自 HW1、HW4、HW6、HW7：
\begin{itemize}
  \item \textbf{HW4}：STM32 作為 BLE GATT server，把加速度計數值 notify 給 client。這次改成只傳「打擊事件」，並同時連兩塊板子。
  \item \textbf{HW7}：CMSIS-DSP 數位濾波。這次濾波後要即時做峰值偵測與分類，而不只是顯示波形。
  \item \textbf{HW1、HW6}：timer 中斷觸發取樣、semaphore 與多個 RTOS task 分工。這次取樣率提高到 400 Hz，對時間精準度要求更高。
\end{itemize}

歷屆作品「空氣爵士鼓」證明了用 STM32 偵測敲擊可行，但它是自由演奏，打到就發聲，無法衡量打得準不準。
本專題的不同在於：
\begin{enumerate}
  \item \textbf{有遊戲畫面與譜面}：我們自製節奏遊戲，包含音符捲動、判定、連擊與計分。
  \item \textbf{時間判定}：節奏遊戲對延遲非常敏感，必須量化並壓低「敲擊 → 判定」的端到端延遲，並提供延遲校正。
  \item \textbf{咚／咔分類與大音符}：依揮動方向分出咚、咔，兩手同時敲擊則判定為大音符，需要對齊兩塊板子的事件時間。
\end{enumerate}
這三點分別對應課程要求的多執行緒 RTOS 設計、感測資料處理與傳輸，以及效能最佳化。

\section{預定作法}

系統分為三段：兩塊 STM32 各自偵測敲擊、BLE 只傳打擊事件、電腦端遊戲依時間判定並顯示。

\begin{figure}[H]
\centering
\resizebox{\linewidth}{!}{%
\begin{tikzpicture}[
  font=\small,
  blk/.style={draw, rounded corners=3pt, minimum width=2.6cm, minimum height=1.25cm, align=center, fill=white},
  hot/.style={blk, draw=blue!60!black, fill=blue!8, line width=0.9pt},
  arr/.style={-{Stealth[length=2.2mm]}, thick},
  lbl/.style={font=\footnotesize, text=black!65}
]
% left stick
\node[blk] (s1) {\textbf{LSM6DSL}\\[-1pt]{\footnotesize 400 Hz 取樣}};
\node[hot, right=1.0cm of s1] (d1) {\textbf{Task\_DSP}\\[-1pt]{\footnotesize 偵測、咚咔分類}};
\node[blk, right=1.0cm of d1] (b1) {\textbf{Task\_BLE}\\[-1pt]{\footnotesize notify 事件}};
% right stick
\node[blk, below=1.5cm of s1] (s2) {\textbf{LSM6DSL}\\[-1pt]{\footnotesize 400 Hz 取樣}};
\node[hot, right=1.0cm of s2] (d2) {\textbf{Task\_DSP}\\[-1pt]{\footnotesize 偵測、咚咔分類}};
\node[blk, right=1.0cm of d2] (b2) {\textbf{Task\_BLE}\\[-1pt]{\footnotesize notify 事件}};
% PC
\node[blk, minimum width=3.4cm, minimum height=3.9cm, right=2.0cm of $(b1)!0.5!(b2)$, anchor=west] (pc)
  {\textbf{電腦（pygame）}\\[4pt]{\footnotesize 合併兩支鼓棒事件}\\{\footnotesize 大音符判定}\\{\footnotesize 良／可／不可}\\{\footnotesize 譜面捲動、計分}};
\begin{scope}[on background layer]
  \node[draw=black!40, fill=black!4, rounded corners=4pt, inner sep=8pt, fit=(s1)(d1)(b1),
        label={[anchor=south west, font=\small\bfseries]north west:左鼓棒：STM32 \#1（FreeRTOS）}] (L) {};
  \node[draw=black!40, fill=black!4, rounded corners=4pt, inner sep=8pt, fit=(s2)(d2)(b2),
        label={[anchor=south west, font=\small\bfseries]north west:右鼓棒：STM32 \#2（FreeRTOS）}] (R) {};
\end{scope}
\draw[arr] (s1) -- node[lbl, above]{queue} (d1);
\draw[arr] (d1) -- node[lbl, above]{事件} (b1);
\draw[arr] (s2) -- node[lbl, above]{queue} (d2);
\draw[arr] (d2) -- node[lbl, above]{事件} (b2);
\draw[arr] (b1.east) -- node[lbl, above, sloped]{BLE} (pc.west |- b1.east);
\draw[arr] (b2.east) -- node[lbl, below, sloped]{BLE} (pc.west |- b2.east);
\end{tikzpicture}}
\caption{圖 1　系統架構：敲擊在各自的 STM32 上偵測與分類，BLE 只傳打擊事件，電腦端合併兩支鼓棒並依時間判定。}
\end{figure}

\subsection{STM32 上的 RTOS 任務（兩塊板子相同韌體）}

\noindent{\small\begin{tabularx}{\linewidth}{@{}>{\raggedright\arraybackslash}p{2.6cm} Y Y l@{}}
\toprule
Task（優先權） & 觸發方式 & 工作內容 & 延伸自 \\
\midrule
Task\_Sensor（最高） & Timer 中斷每 2.5 ms 給 semaphore（400 Hz） & 讀 LSM6DSL 加速度計＋陀螺儀共 6 軸，加上時間戳放入 queue & HW1、6 \\
Task\_DSP（中） & 從 queue 收資料 & 濾波、峰值偵測、去彈跳、咚／咔分類、估計力道 & HW7 \\
Task\_BLE（低） & 打擊事件或 GATT 事件 & 連線／斷線、notify 打擊事件、接收設定寫入 & HW4 \\
\bottomrule
\end{tabularx}}

開機時以 user button 設定這塊板子是左手或右手，並用 LED 閃爍回饋每次偵測到的敲擊（延伸自 HW1）。

\subsection{敲擊偵測與咚／咔分類}
\begin{enumerate}
  \item \textbf{濾波}：以 CMSIS-DSP FIR 濾除重力等低頻成分與高頻雜訊，保留敲擊的衝擊訊號。
  \item \textbf{峰值偵測}：加速度大小超過門檻且為局部最大值時視為一次敲擊；之後約 80 ms 內不重複觸發，避免一次敲擊被算成兩次。
  \item \textbf{咚／咔分類}：依敲擊瞬間衝擊的主要方向判斷——往下敲為「咚」，往外側揮為「咔」。先用門檻規則實作，
        再錄下每人每種敲法數十次，檢查準確率並調整門檻；必要時改用簡單的決策樹。
  \item \textbf{力道}：以峰值大小估計，作為遊戲中的音量與特效依據。
\end{enumerate}

\subsection{BLE GATT 設計}

\noindent{\small\begin{tabularx}{\linewidth}{@{}l l Y@{}}
\toprule
Characteristic & 屬性 & 內容 \\
\midrule
Hit & Notify & 鼓棒編號、咚／咔、力道、板上時間戳 \\
Raw Data & Notify & 6 軸原始資料，只在調校模式開啟，用來錄資料與調門檻 \\
Config & Write & 偵測門檻、取樣頻率、模式（延續 HW4 由 client 寫入設定的做法） \\
\bottomrule
\end{tabularx}}

為降低延遲，連線時要求最短的 connection interval（7.5 ms），讓事件能盡快送出。

\subsection{電腦端遊戲}
\begin{itemize}
  \item \textbf{畫面}：以 pygame 自製，包含音符從右往左捲動、判定區、判定文字、連擊數與分數。
  \item \textbf{譜面}：以簡單的文字檔記錄歌曲 BPM 與每個音符的拍點與種類（咚、咔、大咚、大咔），至少做 2 首歌。
  \item \textbf{判定}：打擊時間與音符時間相差 $\pm$50 ms 內為「良」、$\pm$100 ms 內為「可」，其餘為「不可」（窗口可調）。
        兩支鼓棒在 40 ms 內同時打同一種，判定為大音符。
  \item \textbf{延遲校正}：遊戲提供校正模式，玩家跟著節拍器敲擊，自動估計系統延遲並套用補償。
\end{itemize}

\subsection{量測與最佳化}
\begin{itemize}
  \item \textbf{板上偵測延遲}：偵測到敲擊時拉高 GPIO，用邏輯分析儀量測 Task\_DSP 的處理時間（延伸自 HW5）。
  \item \textbf{端到端延遲}：用手機 240 fps 慢動作同時拍攝敲擊動作與螢幕，計算「敲擊 → 畫面判定」的延遲分布。
  \item \textbf{定點運算}：FIR 分別以 float32 與 q15 實作，用 DWT cycle counter 比較執行時間與誤差。
  \item \textbf{傳輸量}：只傳打擊事件，相較於連續串流原始資料（400 Hz × 6 軸 × 2 bytes ≈ 4.8 KB/s）大幅降低。
\end{itemize}

\section{預期結果與時程}

期末 demo 時兩手拿著 STM32 打完一首歌，畫面即時顯示判定與分數，並以數據說明偵測準確率與延遲。

\noindent{\small\begin{tabularx}{\linewidth}{@{}l Y@{}}
\toprule
指標 & 目標 \\
\midrule
敲擊偵測 & 漏打與誤觸各 $<$ 5\% \\
咚／咔分類 & 三位組員各敲每種 30 次，準確率 $\geq$ 90\% \\
大音符 & 兩手同時敲擊的判定成功率 $\geq$ 85\% \\
端到端延遲 & 敲擊到畫面判定 $<$ 50 ms（校正前量測，校正後補償） \\
遊戲內容 & 自製遊戲畫面＋至少 2 首譜面，可完整玩完一首 \\
定點加速 & 量出 q15 相對 float32 的 FIR 執行時間與誤差 \\
\bottomrule
\end{tabularx}}

\subsection*{時程}
\noindent{\small\begin{tabularx}{\linewidth}{@{}l Y@{}}
\toprule
週次 & 工作 \\
\midrule
Week 11 & Proposal 報告；以 HW4 的 BLE GATT server 為基礎建立專案；遊戲畫面雛形 \\
Week 12 & 單塊板子 400 Hz 取樣與敲擊偵測；電腦端收到事件並在畫面上顯示 \\
Week 13 & 咚／咔分類；兩塊板子同時連線；譜面與判定可玩（Progress 報告） \\
Week 14 & 大音符、延遲校正模式；q15 與 float32 比較 \\
Week 15 & 準確率與延遲量測、第二首譜面、整理 demo \\
\bottomrule
\end{tabularx}}

\subsection*{分工}
\begin{itemize}
  \item 李健誌：STM32 RTOS 架構、BLE GATT server 與雙板連線
  \item 黃琮堡：敲擊偵測、咚／咔分類與 DSP 最佳化
  \item 莊景皓：pygame 遊戲畫面、譜面與判定系統
\end{itemize}
（分工為暫定，可依組員意願調整；延遲量測三人一起做。）

\section{參考資料}
\begin{enumerate}[label={[\arabic*]}]
  \item 課程講義：CMSIS-RTOS2、STM32 IoT Node Onboard Sensors、BLE Introduction Labs、Data Acquisition、CMSIS-DSP CubeIDE、STM32 Timer and PWM。
  \item 本組 HW1（RTOS API）、HW4（BLE GATT server）、HW5（PWM 與邏輯分析儀）、HW6（Timer＋ADC＋DMA）、HW7（CMSIS-DSP）程式碼。
  \item 嵌入式系統專題簡介（歷屆期末展示：空氣爵士鼓）。
  \item STMicroelectronics, \href{https://www.st.com/en/evaluation-tools/b-l475e-iot01a.html}{B-L475E-IOT01A Discovery kit for IoT node}.
  \item STMicroelectronics, \href{https://www.st.com/resource/en/datasheet/lsm6dsl.pdf}{LSM6DSL datasheet}.
  \item Arm, \href{https://github.com/ARM-software/CMSIS-DSP}{CMSIS-DSP}.
  \item \href{https://github.com/hbldh/bleak}{bleak: Bluetooth Low Energy platform agnostic client for Python}.
  \item \href{https://www.pygame.org/docs/}{pygame documentation}.
  \item Bandai Namco,《太鼓之達人》（遊戲玩法與判定概念參考）。
\end{enumerate}

\end{document}
